\chapter{Permutation core}
\label{ch:ArchonFull600Theory_Basic}

\section{Complete permutation actions}

Let \(\Sigma\) be the finite set of all sticker slots and let \(G\) be
the legal subgroup of its permutation group. Words act chronologically:
\(P Q\) acts by \(Q\circ P\). A cycle \((a\,b\,c)\) sends \(a\) to \(b\),
\(b\) to \(c\), and \(c\) to \(a\), fixing the other slots. Write
\(\operatorname{supp}(P)=\{s\in\Sigma:P(s)\ne s\}\) for the full support.

\begin{lemma}[Agreement on a target orbit]
  \label{lem:scope-same-restriction}
  \lean{ArchonFull600Theory.scopeSameRestriction}
  \leanok
  % SOURCE: Theory supplement, Scope tests, lines 1528--1533 (read from references/Full_600cell_Puzzle_Theory.tex)
  % SOURCE QUOTE: First, on six stickers the permutations \begin{equation}
  % u=(1\,2\,3),
  % \qquad
  % v=(1\,2\,3)(4\,5\,6)
  % \end{equation} have identical restrictions to the target orbit
  % \(\{1,2,3\}\) but different complete actions. Both are even.
  \textit{Source: Theory supplement, Scope test S1.}
  On the six stickers \(1,\ldots,6\), let \(u=(1\,2\,3)\) and
  \(v=(1\,2\,3)(4\,5\,6)\). Both are even, their restrictions to
  \(\{1,2,3\}\) agree, and \(u\ne v\).
\end{lemma}
% SOURCE QUOTE PROOF: \end{equation} have identical restrictions to the target orbit
% \(\{1,2,3\}\) but different complete actions. Both are even.
\begin{proof}
  Each three-cycle is a product of two transpositions. The extra cycle
  in \(v\) fixes all three target slots, and \(u(4)=4\) while \(v(4)=5\).
\end{proof}

\begin{lemma}[Disjoint visible supports]
  \label{lem:scope-visible-disjoint}
  \lean{ArchonFull600Theory.scopeVisibleDisjoint}
  \leanok
  % SOURCE: Theory supplement, Scope tests, lines 1539--1546 (read from references/Full_600cell_Puzzle_Theory.tex)
  % SOURCE QUOTE: Second, consider \begin{equation}
  % a=(1\,2)(5\,6),
  % \qquad
  % b=(3\,4)(6\,7).
  % \end{equation} If only stickers \(1,2,3,4\) are displayed, their visible
  % supports are disjoint. But their full actions do not commute: applying
  % \(a\) then \(b\) sends \(5\) to \(7\), whereas applying \(b\) then \(a\)
  % sends \(5\) to \(6\).
  \textit{Source: Theory supplement, Scope test S2.}
  On the seven stickers \(1,\ldots,7\), let \(a=(1\,2)(5\,6)\) and
  \(b=(3\,4)(6\,7)\). Their supports intersect the displayed set
  \(\{1,2,3,4\}\) in the disjoint sets \(\{1,2\}\) and \(\{3,4\}\).
  The full actions fail to commute, with \(b(a(5))=7\) and \(a(b(5))=6\).
\end{lemma}
% SOURCE QUOTE PROOF: applying
% \(a\) then \(b\) sends \(5\) to \(7\), whereas applying \(b\) then \(a\)
% sends \(5\) to \(6\).
\begin{proof}
  The transpositions give the displayed intersections by inspection.
  The two action orders send \(5\) along \(5\mapsto6\mapsto7\) and
  \(5\mapsto5\mapsto6\), respectively.
\end{proof}

\begin{lemma}[Commutation on disjoint full supports]
  \label{lem:disjoint-full-supports}
  \lean{ArchonFull600Theory.disjointFullSupportsCommute}
  \leanok
  % SOURCE: Theory supplement, Scope tests, lines 1548--1551 (read from references/Full_600cell_Puzzle_Theory.tex)
  % SOURCE QUOTE: permutations with disjoint full supports commute.
  \textit{Source: Theory supplement, Scope test S3.}
  If \(\operatorname{supp}(P)\) and \(\operatorname{supp}(Q)\) are
  disjoint, then \(Q\circ P=P\circ Q\) on \(\Sigma\).
\end{lemma}
% SOURCE QUOTE PROOF: Each
% permutation fixes the support of the other pointwise, and a permutation
% maps its own support to itself; checking a point in either support or
% outside both gives equality of the two compositions.
\begin{proof}
  A permutation preserves its full support by injectivity. If \(s\) is
  in the support of \(P\), then \(s\) and \(P(s)\) are fixed by \(Q\), so
  both compositions send \(s\) to \(P(s)\). The argument interchanging
  \(P\) and \(Q\) covers the support of \(Q\). Both fix every other \(s\).
\end{proof}

\begin{definition}[Complete nonfixed map]
  \label{def:nonfixed-map}
  \lean{ArchonFull600Theory.nonfixedMap}
  \leanok
  % SOURCE: Theory supplement, Lemma 14, lines 1434--1437 (read from references/Full_600cell_Puzzle_Theory.tex)
  % SOURCE QUOTE: E_W=\{(u,W(u)):W(u)\ne u\}.
  \textit{Source: Theory supplement, Lemma~14.}
  For a permutation \(W\), let \(E_W=\{(u,W(u)):W(u)\ne u\}\) be its
  complete nonfixed map, a set of ordered pairs in \(\Sigma\times\Sigma\).
\end{definition}
\begin{proof}
  Each moved source has exactly its image as a target. Reconstruct \(W\)
  by fixing every slot absent from the first coordinates of \(E_W\).
\end{proof}

\begin{lemma}[Complete endpoint transport]
  \label{lem:endpoint-transport}
  \lean{ArchonFull600Theory.transportNonfixedMap}
  \leanok
  \uses{def:nonfixed-map}
  % SOURCE: Theory supplement, Lemma 14, lines 1433--1444 (read from references/Full_600cell_Puzzle_Theory.tex)
  % SOURCE QUOTE: Let \(W\) and \(S\) be verified legal sticker permutations, with
  % chronological witnesses. Suppose the complete nonfixed map of \(W\) is
  % \begin{equation}
  % E_W=\{(u,W(u)):W(u)\ne u\}.
  % \end{equation} Then the complete nonfixed map of the chronological
  % conjugate \(S^{-1}WS\) is \begin{equation}
  % E_{S^{-1}WS}
  % =
  % \{(S(u),S(W(u))):(u,W(u))\in E_W\}.
  % \end{equation} Thus transporting both endpoints computes exactly the
  % same full permutation as replaying the conjugated legal witness,
  % including all collateral.
  \textit{Source: Theory supplement, Lemma~14.}
  Let \(W,S\in G\) have chronological legal witnesses. The complete
  nonfixed map of the chronological conjugate satisfies
  \[
    E_{S^{-1}WS}=\{(S(u),S(v)):(u,v)\in E_W\}.
  \]
\end{lemma}
% SOURCE QUOTE PROOF: As an ordinary function, the chronological word \(S^{-1}WS\) is
% \begin{equation}
% S\circ W\circ S^{-1}.
% \end{equation} For each \(u\), \begin{equation}
% (S\circ W\circ S^{-1})(S(u))=S(W(u)).
% \end{equation} Since \(S\) is bijective, \(S(W(u))\ne S(u)\) exactly
% when \(W(u)\ne u\). Every moved source of the conjugate is therefore
% represented, and no fixed source is included.
\begin{proof}
  The action at \(x\) is \(S(W(S^{-1}(x)))\). Write \(x=S(u)\); its
  target is \(S(W(u))\), which differs from \(S(u)\) exactly when
  \(W(u)\ne u\). Every \(x\) has a unique such \(u\). The pair-set
  equality determines the conjugated action on all slots, including
  collateral cycles.
\end{proof}

\section{Independent alternating actions}

\begin{definition}[Pure-controller certificate]
  \label{def:pure-controller}
  \lean{ArchonFull600Theory.PureControllerCertificate}
  \leanok
  % SOURCE: Theory supplement, hypothesis C, lines 851--860 (read from references/Full_600cell_Puzzle_Theory.tex)
  % SOURCE QUOTE: \textbf{C --- Nine pure-controller certificates.} For \begin{equation}
  % \mathcal I=\{1,4,5,7,18,19,27,31,33\},
  % \end{equation} each orbit has a legal word whose complete sticker
  % permutation is a single three-cycle in that orbit and fixes every other
  % sticker. A verified legal relocation produces a cycle
  % \((A_o\,B_o\,C_o)\). For every other position \(X\) in that orbit, a
  % verified legal setup fixes the two buffer sticker slots \(A_o,B_o\)
  % pointwise and sends \(C_o\) to \(X\). Coverage means valid transitions
  % and coverage of every nonbuffer destination, not merely a reported
  % search-node count.
  \textit{Source: Theory supplement, hypothesis C.}
  A certificate for a selected orbit \(O\subseteq\Sigma\) gives a legal
  seed whose complete permutation is \((a\,b\,c)\) for distinct
  \(a,b,c\in O\), and a legal relocation sending \(a,b,c\) respectively
  to the distinct slots \(A,B,C\in O\). For each \(X\in O\setminus\{A,B\}\)
  there is a legal setup \(S_X\) fixing \(A,B\) individually and sending
  \(C\) to \(X\). The seed, relocation, and setups belong to \(G\).
\end{definition}
\begin{proof}
  The seed condition specifies the complete permutation, with the
  identity at each slot outside \(\{a,b,c\}\). The coverage condition
  includes \(C\), where the identity serves as a setup.
\end{proof}

\begin{lemma}[Independent supported alternating actions]
  \label{lem:supported-alternating-actions}
  \lean{ArchonFull600Theory.supportedAlternatingActions}
  \leanok
  \uses{def:pure-controller,lem:endpoint-transport,lem:disjoint-full-supports}
  % SOURCE: Theory supplement, Lemma 9, lines 1087--1090 (read from references/Full_600cell_Puzzle_Theory.tex)
  % SOURCE QUOTE: Under R and C, \(G\) contains a subgroup acting as \begin{equation}
  % A_{3600}\times A_{2400}\times A_{7200}^{\,7}
  % \end{equation} on the nine selected single-sticker orbits and fixing
  % every other sticker.
  \textit{Source: Theory supplement, Lemma~9.}
  Let \(O_i\subseteq\Sigma\) be the pairwise disjoint selected orbits for
  \(i\in\mathcal I=\{1,4,5,7,18,19,27,31,33\}\), preserved setwise by
  \(G\). Suppose \(O_1\) has size \(3600\), \(O_{33}\) has size \(2400\),
  the other seven have size \(7200\), and each has a pure-controller
  certificate. There is an injective homomorphism
  \[
    \prod_{i\in\mathcal I}\operatorname{Alt}(O_i)\longrightarrow G
  \]
  sending a family of even permutations to its prescribed action on
  each \(O_i\) and fixing every other sticker. The degrees give
  \(A_{3600}\times A_{2400}\times A_{7200}^{\,7}\).
\end{lemma}
% SOURCE QUOTE PROOF: Words execute from left to right. Let \begin{equation}
% t_C=(A\,B\,C)
% \end{equation} be a relocated pure controller. If \(S\) fixes \(A,B\)
% and sends \(C\) to \(X\), then the chronological word \begin{equation}
% S^{-1}t_CS
% \end{equation} has ordinary functional action
% \(S\circ t_C\circ S^{-1}\), so it is exactly \begin{equation}
% t_X=(A\,B\,X).
% \end{equation} Conjugating the complete pure permutation gives a
% complete pure permutation. Temporary collateral during its legal witness
% is irrelevant to this net support statement.
%
% For distinct nonbuffer positions \(X,Y\), direct evaluation gives
% \begin{equation}
% t_Y^{-1}t_X=(A\,Y\,X).
% \end{equation} Together with the original stars and their inverses,
% these supply every three-cycle containing \(A\). For distinct
% \(x,y,z\ne A\), \begin{equation}
% (A\,x\,y)(A\,z\,x)=(x\,y\,z).
% \end{equation} Thus the stars generate every three-cycle.
%
% Every even permutation is a product of three-cycles. For example,
% express it as an even number of transpositions and pair them; a pair
% sharing a point is a three-cycle, identical transpositions cancel, and a
% disjoint pair can be rewritten as \begin{equation}
% (a\,b)(c\,d)=(a\,b\,c)(a\,d\,c)
% \end{equation} in the chronological convention. Conversely, every star
% is even. The supported subgroup generated on this orbit is therefore
% exactly \(A_N\).
%
% The nine supported subgroups have disjoint full sticker supports. They
% commute, and the multiplication map from their direct product is
% injective: restricting an identity product to one selected orbit makes
% that factor the identity. This proves the claim.
\begin{proof}
  Transport the seed along the legal relocation to obtain the full
  cycle \(t_C=(A\,B\,C)\). Transport \(t_C\) along each setup to obtain
  \(t_X=(A\,B\,X)\). The conjugated witnesses are legal by subgroup
  closure, and each resulting cycle fixes the entire complement of
  its three slots.

  For distinct nonbuffer \(X,Y\), chronological evaluation gives
  \(t_Y^{-1}t_X=(A\,Y\,X)\): the product sends \(A\) to \(Y\), \(Y\) to
  \(X\), \(X\) to \(A\), and fixes \(B\) and all other slots. Together
  with the stars and their inverses this gives every three-cycle
  containing \(A\). For distinct \(x,y,z\ne A\), the chronological
  product \((A\,x\,y)(A\,z\,x)\) fixes \(A\) and cycles \(x,y,z\).
  Thus every three-cycle on \(O_i\), extended by the identity, is legal.

  An even permutation is a product of pairs of transpositions. Equal
  transpositions cancel, a pair sharing one point forms a three-cycle,
  and a disjoint pair factors as
  \((a\,b)(c\,d)=(a\,b\,c)(a\,d\,c)\) in chronological order.
  Hence all even permutations of \(O_i\) extend by the identity to
  elements of \(G\). Conversely, each star restricts to an even
  permutation on \(O_i\).

  The extended factors commute by disjointness of their full supports,
  so their product is a homomorphism from the direct product. Its
  restriction to \(O_i\) recovers the \(i\)-th factor, proving
  injectivity and the prescribed action. The factors each fix the
  complement of the union of the orbits.
\end{proof}

\section{Reachable fixed-frame colorings}

\begin{definition}[Reachable colorings]
  \label{def:reachable-colorings}
  \lean{ArchonFull600Theory.reachableColorings}
  \leanok
  % SOURCE: Theory supplement, Lemma 10, lines 1138--1141 and 1176--1178 (read from references/Full_600cell_Puzzle_Theory.tex)
  % SOURCE QUOTE: the set \(\Omega\) of reachable fixed-frame face-color
  % states
  \textit{Source: Theory supplement, Lemma~10.}
  For a finite palette \(\Lambda\) and solved coloring
  \(c_0\colon\Sigma\to\Lambda\), set
  \(\Omega=\{c_0\circ g^{-1}:g\in G\}\). Positions and symbols are
  read in the fixed spatial frame.
\end{definition}
\begin{proof}
  The color at \(s\) after \(g\) is the solved color of the slot sent
  to \(s\) by \(g\). The identity realizes \(c_0\), and \(\Omega\) is
  finite as a set of maps between finite sets.
\end{proof}

\begin{definition}[Balanced-color count]
  \label{def:balanced-color-count}
  \lean{ArchonFull600Theory.balancedColorCount}
  \leanok
  % SOURCE: Theory supplement, Lemma 10, lines 1142--1145 (read from references/Full_600cell_Puzzle_Theory.tex)
  % SOURCE QUOTE: B=
  % \frac{3600!}{(6!)^{600}}\,
  % \frac{2400!}{(4!)^{600}}\,
  % \left[\frac{7200!}{(12!)^{600}}\right]^7.
  \textit{Source: Theory supplement, Lemma~10.}
  Define the natural number
  \[
    B=\frac{3600!}{(6!)^{600}}\,
      \frac{2400!}{(4!)^{600}}\,
      \left(\frac{7200!}{(12!)^{600}}\right)^7.
  \]
\end{definition}
\begin{proof}
  For \(m=6,4,12\), ordering the \(600m\) slots and taking consecutive
  blocks of size \(m\) partitions the slots into \(600\) labeled
  fibers. Exactly \((m!)^{600}\) orderings yield each such partition,
  so each factorial quotient is exact.
\end{proof}

\begin{lemma}[Balanced-color lower bound]
  \label{lem:balanced-color-lower-bound}
  \lean{ArchonFull600Theory.balancedColorLowerBound}
  \leanok
  \uses{def:pure-controller,def:reachable-colorings,def:balanced-color-count,lem:supported-alternating-actions}
  % SOURCE: Theory supplement, hypothesis C, lines 862--864, and Lemma 10, lines 1138--1147 (read from references/Full_600cell_Puzzle_Theory.tex)
  % SOURCE QUOTE: The solved coloring has \(600\) distinct home-cell symbols, with six
  % copies of each in O01, four in O33, and twelve in each of the other
  % seven selected orbits.
  % SOURCE QUOTE: Under R and C, the set \(\Omega\) of reachable fixed-frame face-color
  % states satisfies \begin{equation}
  % |\Omega|\ge B,
  % \end{equation} where \begin{equation}
  % B=
  % \frac{3600!}{(6!)^{600}}\,
  % \frac{2400!}{(4!)^{600}}\,
  % \left[\frac{7200!}{(12!)^{600}}\right]^7.
  % \end{equation} These states can be attained with every unselected orbit
  % solved at macro boundaries.
  \textit{Source: Theory supplement, Lemma~10 and hypothesis C.}
  Assume the nine orbits and certificates of
  Lemma~\ref{lem:supported-alternating-actions}. Let \(\Lambda\) have
  \(600\) distinct symbols and \(c_0\) have six copies of each on
  \(O_1\), four on \(O_{33}\), and twelve on each of the other seven.
  Then at least \(B\) colorings are attained by legal permutations
  fixing every slot outside the selected orbits. In particular
  \(|\Omega|\ge B\), and each unselected orbit is solved at the
  boundaries of the witnessing macros.
\end{lemma}
% SOURCE QUOTE PROOF: Consider one selected orbit with \(N=600m\) positions and \(m\) copies
% of each color. The stabilizer of its solved coloring in \(S_N\) is
% \begin{equation}
% Q=(S_m)^{600}.
% \end{equation} Since \(m\in\{4,6,12\}\), this stabilizer contains a
% transposition of two equal-colored positions and hence an odd
% permutation. Therefore \begin{equation}
% |A_N\cap Q|=\frac{|Q|}{2}.
% \end{equation} The number of colorings in the \(A_N\) orbit is
% \begin{equation}
% \frac{|A_N|}{|A_N\cap Q|}
% =
% \frac{N!/2}{(m!)^{600}/2}
% =
% \frac{N!}{(m!)^{600}}.
% \end{equation}
%
% The odd same-color transposition is used only to describe the coloring
% stabilizer. It need not itself be a legal move.
%
% Lemma 9 gives independent alternating actions with full purity, so the
% nine factors multiply and every unselected sticker remains solved at the
% ends of the resulting macros. This gives \(B\).
\begin{proof}
  On an orbit of size \(600m\), any two colorings with those fiber sizes
  are related by a permutation obtained by matching their fibers.
  The stabilizer of the solved coloring has order \((m!)^{600}\):
  it permutes each fiber independently. Since \(m\ge 2\), swapping two
  slots of a single fiber gives an odd element of the stabilizer.
  Precomposing an odd matching permutation with this swap makes it
  even and preserves its action on the solved coloring. Thus the
  alternating action attains the full symmetric orbit of size
  \((600m)!/(m!)^{600}\) by orbit--stabilizer.

  Use the supported actions independently for the nine orbits. Different
  choices of orbit colorings give different full colorings, and each
  supported macro fixes the slots outside these orbits at its boundary.
  Multiplying the nine cardinalities gives \(B\).
\end{proof}

\section{Word length and total variation}

For the concrete bounds let the alphabet have the \(1800\) labels
\(\{H_c,T_c,T_c^{-1}:c\in\mathcal C\}\), one for each of the three
letter types at each of \(600\) caps. Each label contributes one to
length. Its interpretation is the corresponding permutation in \(G\).
The legal alphabet realizes \(G\) by words, as its inverses are also
available as letters.

\begin{definition}[Coloring after a word]
  \label{def:word-coloring}
  \lean{ArchonFull600Theory.wordColoring}
  \leanok
  \uses{def:reachable-colorings}
  % SOURCE: Theory supplement, Lemma 11, lines 1191--1194 and 1206--1207 (read from references/Full_600cell_Puzzle_Theory.tex)
  % SOURCE QUOTE: At most \(q^j\) words have length \(j\). Different words may give the
  % same state
  \textit{Source: Theory supplement, Lemma~11.}
  Interpret a word \(w\) as the chronological product \(g_w\) of its
  letters' permutations, starting with the identity. Its output from
  solved is \(x_w=c_0\circ g_w^{-1}\in\Omega\).
\end{definition}
\begin{proof}
  Appending a letter \(a\) composes \(g_w\) on the left by the
  permutation for \(a\). Subgroup closure puts every \(g_w\) in \(G\).
\end{proof}

\begin{definition}[Distance from solved]
  \label{def:word-distance}
  \lean{ArchonFull600Theory.wordDistance}
  \leanok
  \uses{def:word-coloring,def:reachable-colorings}
  % SOURCE: Theory supplement, Lemma 11, lines 1191--1195 (read from references/Full_600cell_Puzzle_Theory.tex)
  % SOURCE QUOTE: Under R and C, use the unit-cost alphabet \begin{equation}
  % \mathcal A=\{H_c,T_c,T_c^{-1}:c\in\mathcal C\},
  % \qquad q=1800.
  % \end{equation} Some reachable face-color state has distance at least
  % \(46648\) from solved.
  \textit{Source: Theory supplement, Lemma~11.}
  For \(x\in\Omega\), define \(d(c_0,x)\) as the least length of a
  word \(w\) with \(x_w=x\), using a finite alphabet realizing \(G\).
\end{definition}
\begin{proof}
  Since a legal permutation attains \(x\) and the alphabet realizes
  \(G\), the set of such lengths is nonempty and has a least element.
\end{proof}

\begin{lemma}[Word-ball count]
  \label{lem:word-ball-count}
  \lean{ArchonFull600Theory.wordBallCount}
  \leanok
  \uses{def:reachable-colorings,def:word-distance,def:word-coloring}
  % SOURCE: Theory supplement, Lemma 11, lines 1206--1213 (read from references/Full_600cell_Puzzle_Theory.tex)
  % SOURCE QUOTE: |\{x:d(\mathrm{solved},x)\le d\}|
  % \le
  % \sum_{j=0}^dq^j
  % =
  % \frac{q^{d+1}-1}{q-1}.
  \textit{Source: Theory supplement, Lemma~11.}
  If the finite alphabet realizing \(G\) has \(q\ge 2\) letters, then
  for each \(r\ge 0\),
  \[
    |\{x\in\Omega:d(c_0,x)\le r\}|\le
    \sum_{j=0}^{r}q^j=\frac{q^{r+1}-1}{q-1}.
  \]
\end{lemma}
% SOURCE QUOTE PROOF: At most \(q^j\) words have length \(j\). Different words may give the
% same state
\begin{proof}
  Words of length \(j\) are \(j\)-tuples of letters, so there are \(q^j\)
  such words. The state ball is the image of the disjoint union of
  these word sets over \(0\le j\le r\). The geometric sum identity
  follows by multiplication by \(q-1>0\) and telescoping.
\end{proof}

For the integer comparisons below put \(K=10^{151851}\). Evaluate
factorials and powers by successive integer products. In the
calculation of \(B\), the common denominator and numerator are
\[
  D=(6!)^{600}(4!)^{600}(12!)^{4200},\qquad
  P=3600!\,2400!\,(7200!)^7.
\]
The three exact divisions defining \(B\) give \(D B=P\). Each floor
\(\lfloor n/d\rfloor=k\) below can be checked by the two integer
comparisons \(kd\le n<(k+1)d\).

\begin{lemma}[Scale of the color count]
  \label{lem:n7-color-count}
  \lean{ArchonFull600Theory.balancedColorCountGtPower}
  \leanok
  \uses{def:balanced-color-count}
  % SOURCE: Theory supplement, Distinct symmetry and motion groups, lines 1358--1360 (read from references/Full_600cell_Puzzle_Theory.tex)
  % SOURCE QUOTE: Under C, the
  % lower bound \(B>10^{151851}\) alone already separates the scale of the
  % legal action from the finite geometric group.
  \textit{Source: Theory supplement, Distinct symmetry and motion groups, and handoff check N7.}
  In the natural numbers, \(B>10^{151851}\).
\end{lemma}
\begin{proof}
  Exact evaluation gives \(\lfloor 1000P/(DK)\rfloor=1459\), so
  \(1000B\ge1459K>1000K\).
\end{proof}

\begin{lemma}[Radius \(46647\) comparison]
  \label{lem:n8-radius-lower}
  \lean{ArchonFull600Theory.radius46647Comparison}
  \leanok
  \uses{def:balanced-color-count,lem:n7-color-count}
  % SOURCE: Theory supplement, Lemma 11, lines 1215--1218 (read from references/Full_600cell_Puzzle_Theory.tex)
  % SOURCE QUOTE: 1800^{46648}-1<1799B,
  \textit{Source: Theory supplement, Lemma~11, and handoff check N8.}
  In the natural numbers, \(1800^{46648}-1<1799B\).
\end{lemma}
% SOURCE QUOTE PROOF: They can be evaluated without logarithms: form the
% factorials in \(B\) as integer products, perform the indicated exact
% multinomial divisions, and compare positive integers.
\begin{proof}
  Exact evaluation gives
  \(\lfloor1000(1800^{46648}-1)/(1799K)\rfloor=4\).
  Thus \(1800^{46648}-1<1799K<1799B\).
\end{proof}

\begin{lemma}[Radius \(46648\) cross-check]
  \label{lem:n8-radius-cross-check}
  \lean{ArchonFull600Theory.radius46648CrossCheck}
  \leanok
  \uses{def:balanced-color-count}
  % SOURCE: Theory supplement, Lemma 11, lines 1225--1229 (read from references/Full_600cell_Puzzle_Theory.tex)
  % SOURCE QUOTE: As an arithmetic
  % cross-check, \begin{equation}
  % 1800^{46649}-1\ge1799B.
  % \end{equation}
  \textit{Source: Theory supplement, Lemma~11, and handoff check N8.}
  In the natural numbers, \(1800^{46649}-1\ge1799B\).
\end{lemma}
\begin{proof}
  Exact evaluation gives
  \(\lfloor1000(1800^{46649}-1)/(1799K)\rfloor=8954\) and
  \(\lfloor1000P/(DK)\rfloor=1459\). Hence
  \(1800^{46649}-1\ge 2\cdot1799K\) and \(B<2K\).
\end{proof}

By Lemmas~\ref{lem:n8-radius-lower} and~\ref{lem:n8-radius-cross-check}
and monotonicity of the geometric sum, \(46648\) is the first radius
at which the word-count upper bound reaches \(B\).

\begin{lemma}[Thousand-letter support comparison]
  \label{lem:n9-support-comparison}
  \lean{ArchonFull600Theory.thousandLetterSupportComparison}
  \leanok
  \uses{def:balanced-color-count,lem:n7-color-count}
  % SOURCE: Theory supplement, Lemma 11, lines 1220--1221 (read from references/Full_600cell_Puzzle_Theory.tex)
  % SOURCE QUOTE: 1800^{1000}\,10^{148595}<B.
  \textit{Source: Theory supplement, Lemma~11, and handoff check N9.}
  In the natural numbers, \(1800^{1000}10^{148595}<B\).
\end{lemma}
% SOURCE QUOTE PROOF: They can be evaluated without logarithms: form the
% factorials in \(B\) as integer products, perform the indicated exact
% multinomial divisions, and compare positive integers.
\begin{proof}
  Exact evaluation gives
  \(\lfloor1000\cdot1800^{1000}10^{148595}/K\rfloor=187\).
  Thus \(1800^{1000}10^{148595}<K<B\).
\end{proof}

\begin{lemma}[Worst-case distance]
  \label{lem:word-distance-lower-bound}
  \lean{ArchonFull600Theory.wordDistanceLowerBound}
  \leanok
  \uses{def:pure-controller,def:reachable-colorings,def:word-distance,lem:balanced-color-lower-bound,lem:word-ball-count,lem:n8-radius-lower}
  % SOURCE: Theory supplement, Lemma 11, lines 1191--1195 (read from references/Full_600cell_Puzzle_Theory.tex)
  % SOURCE QUOTE: Under R and C, use the unit-cost alphabet \begin{equation}
  % \mathcal A=\{H_c,T_c,T_c^{-1}:c\in\mathcal C\},
  % \qquad q=1800.
  % \end{equation} Some reachable face-color state has distance at least
  % \(46648\) from solved.
  \textit{Source: Theory supplement, Lemma~11.}
  Assume the orbit, certificate, and solved-coloring hypotheses of
  Lemma~\ref{lem:balanced-color-lower-bound} and the specified
  \(1800\)-letter unit-cost alphabet realizing \(G\). There is an
  \(x\in\Omega\) with \(d(c_0,x)\ge46648\).
\end{lemma}
% SOURCE QUOTE PROOF: By Equation\textasciitilde{}\eqref{eq:blueprint-1}, the radius-\(46647\)
% ball contains fewer than \(B\) states. Lemma 10 gives at least \(B\)
% reachable states, so some state lies outside that ball.
\begin{proof}
  At radius \(46647\), the word-ball bound is
  \((1800^{46648}-1)/1799<B\) by the strict integer comparison.
  Since \(|\Omega|\ge B\), some reachable state is outside this ball
  and its integer distance is at least \(46648\).
\end{proof}

\begin{definition}[Total variation on a finite set]
  \label{def:total-variation}
  \lean{ArchonFull600Theory.totalVariation}
  \leanok
  % SOURCE: Theory supplement, Lemma 11, lines 1237--1239 (read from references/Full_600cell_Puzzle_Theory.tex)
  % SOURCE QUOTE: \|\mu-\nu\|_{\mathrm{TV}}
  % =\max_E|\mu(E)-\nu(E)|,
  \textit{Source: Theory supplement, Lemma~11.}
  On a finite set, a probability law assigns nonnegative real weights
  summing to one, and \(\mu(E)\) is their sum over the event \(E\).
  Set \(\operatorname{TV}(\mu,\nu)=\max_E|\mu(E)-\nu(E)|\),
  where the maximum ranges over all subsets.
\end{definition}
\begin{proof}
  There are finitely many subsets, including the empty event, so the
  set of event discrepancies has a maximum.
\end{proof}

\begin{lemma}[Thousand-letter total variation]
  \label{lem:total-variation-lower-bound}
  \lean{ArchonFull600Theory.totalVariationLowerBound}
  \leanok
  \uses{def:pure-controller,def:reachable-colorings,def:word-coloring,def:total-variation,lem:balanced-color-lower-bound,lem:n9-support-comparison}
  % SOURCE: Theory supplement, Lemma 11, lines 1197--1202 (read from references/Full_600cell_Puzzle_Theory.tex)
  % SOURCE QUOTE: If \(\mu_{1000}\) is any distribution produced by \(1000\) successive
  % letters of this alphabet from solved and \(\nu\) is uniform on
  % \(\Omega\), then \begin{equation}
  % \|\mu_{1000}-\nu\|_{\mathrm{TV}}
  % >1-10^{-148595}.
  % \end{equation}
  \textit{Source: Theory supplement, Lemma~11.}
  Assume the orbit, certificate, and solved-coloring hypotheses of
  Lemma~\ref{lem:balanced-color-lower-bound} and the specified
  \(1800\)-letter alphabet realizing \(G\). Let \(\mu_{1000}\) be any
  probability law on \(\Omega\) supported on outputs of length-\(1000\)
  words from solved, and let \(\nu\) be uniform on \(\Omega\). Then
  \[
    \operatorname{TV}(\mu_{1000},\nu)>1-10^{-148595}.
  \]
\end{lemma}
% SOURCE QUOTE PROOF: Let \(A\) be the support of \(\mu_{1000}\). Then \begin{equation}
% |A|\le1800^{1000},\qquad \mu_{1000}(A)=1.
% \end{equation} Using the definition \begin{equation}
% \|\mu-\nu\|_{\mathrm{TV}}
% =\max_E|\mu(E)-\nu(E)|,
% \end{equation} we obtain \begin{equation}
% \begin{aligned}
% \|\mu_{1000}-\nu\|_{\mathrm{TV}}
% &\ge1-\frac{|A|}{|\Omega|}\\
% &\ge1-\frac{1800^{1000}}{B}\\
% &>1-10^{-148595},
% \end{aligned}
% \end{equation} where strictness follows from
% Equation\textasciitilde{}\eqref{eq:blueprint-2}.
\begin{proof}
  The support \(A\) is contained in the image of the set of
  length-\(1000\) words, so \(|A|\le1800^{1000}\). All probability mass
  of \(\mu_{1000}\) lies in \(A\), while \(\nu(A)=|A|/|\Omega|\).
  Evaluate the event discrepancy at \(A\) and use \(|\Omega|\ge B\) to
  bound it below by \(1-1800^{1000}/B\). The strict integer
  comparison, divided by the positive numbers \(B\) and \(10^{148595}\),
  makes this greater than \(1-10^{-148595}\).
\end{proof}

\section{Finite group checks}

\begin{lemma}[Derived subgroup of the dihedral group]
  \label{lem:n5-dihedral}
  \lean{ArchonFull600Theory.dihedralFiveCommutatorCard}
  \leanok
  % SOURCE: Theory supplement, Lemma 8, lines 1047--1051 (read from references/Full_600cell_Puzzle_Theory.tex)
  % SOURCE QUOTE: \([D_5,D_5]=\langle R\rangle\).
  \textit{Source: Theory supplement, Lemma~8, and handoff check N5.}
  In the order-ten dihedral group \(D_5\), the derived subgroup
  \([D_5,D_5]\) is the rotation subgroup and has five elements.
\end{lemma}
% SOURCE QUOTE PROOF: For \(D_5\), the quotient by \(\langle R\rangle\) is
% \(C_2\), while \begin{equation}
% [R,F]=R^2
% \end{equation} generates \(\langle R\rangle\). Thus
% \([D_5,D_5]=\langle R\rangle\).
\begin{proof}
  Let \(R(i)=i+1\) and \(F(i)=-i\) on the five residues modulo \(5\).
  The rotation subgroup has order five and abelian quotient of order
  two, so it contains every commutator. The relation
  \(F R F^{-1}=R^{-1}\) gives \([R,F]=R^2\), a generator of that
  subgroup. Both inclusions follow.
\end{proof}

\begin{lemma}[Derived subgroup of the alternating group]
  \label{lem:n5-alternating}
  \lean{ArchonFull600Theory.alternatingFivePerfect}
  \leanok
  % SOURCE: Theory supplement, Lemma 8, lines 1053--1057 (read from references/Full_600cell_Puzzle_Theory.tex)
  % SOURCE QUOTE: For completeness, \(A_5\) is perfect
  \textit{Source: Theory supplement, Lemma~8, and handoff check N5.}
  The derived subgroup of \(A_5=\operatorname{Alt}(\{1,\ldots,5\})\)
  is \(A_5\) itself and has \(60\) elements.
\end{lemma}
% SOURCE QUOTE PROOF: For completeness, \(A_5\) is perfect: take \begin{equation}
% u=(1\,2)(4\,5),\qquad v=(2\,3)(4\,5).
% \end{equation} Their commutator is a three-cycle. The derived subgroup
% is normal, and all three-cycles are conjugate within \(A_5\). Since
% three-cycles generate \(A_5\), its derived subgroup is all of \(A_5\).
\begin{proof}
  The even involutions \(u=(1\,2)(4\,5)\) and \(v=(2\,3)(4\,5)\)
  have a three-cycle as their product, by cancellation of the common
  transposition on the last two slots. Their commutator is a square
  of this three-cycle and is again a three-cycle. Any two three-cycles
  are conjugate by a relabeling, and swapping the two unused labels
  before relabeling corrects its parity if needed. Normality puts all
  three-cycles in the derived subgroup. They generate \(A_5\), of
  order \(5!/2=60\).
\end{proof}

\begin{lemma}[Finite checks of the cycle identities]
  \label{lem:n6-cycle-identities}
  \lean{ArchonFull600Theory.threeCycleIdentityChecks}
  \leanok
  % SOURCE: Theory supplement, Lemma 9, lines 1106--1113 (read from references/Full_600cell_Puzzle_Theory.tex)
  % SOURCE QUOTE: For distinct nonbuffer positions \(X,Y\), direct evaluation gives
  % \begin{equation}
  % t_Y^{-1}t_X=(A\,Y\,X).
  % \end{equation} Together with the original stars and their inverses,
  % these supply every three-cycle containing \(A\). For distinct
  % \(x,y,z\ne A\), \begin{equation}
  % (A\,x\,y)(A\,z\,x)=(x\,y\,z).
  % \end{equation} Thus the stars generate every three-cycle.
  \textit{Source: Theory supplement, Lemma~9, and handoff check N6.}
  In each domain \(\{1,\ldots,n\}\) with \(4\le n\le9\), let
  \(t_Z=(A\,B\,Z)\). For pairwise distinct \(A,B,X,Y\), one has
  \(t_Y^{-1}t_X=(A\,Y\,X)\). For pairwise distinct \(A,x,y,z\), one has
  \((A\,x\,y)(A\,z\,x)=(x\,y\,z)\). Both equalities use chronological
  order and hold on the full domain.
\end{lemma}
% SOURCE QUOTE PROOF: For distinct nonbuffer positions \(X,Y\), direct evaluation gives
% \begin{equation}
% t_Y^{-1}t_X=(A\,Y\,X).
% \end{equation}
\begin{proof}
  Enumerate the ordered distinct quadruples in each domain and compare
  both sides on every slot. In the first comparison the left side
  cycles \(A\mapsto Y\mapsto X\mapsto A\) and fixes \(B\). In the
  second it cycles \(x\mapsto y\mapsto z\mapsto x\) and fixes \(A\).
  Every slot outside the quadruple is fixed in each comparison.
\end{proof}

\section{Starter definition}

\begin{definition}[Starter string]
  \label{def:hello}
  \lean{hello}
  \leanok
  The string named \texttt{hello} is \texttt{world}.
\end{definition}
\begin{proof}
  The string is given directly by its literal.
\end{proof}
